\section{A circle construction with denominator resets}\label{sec:masks}

We now construct the first family of words. A translation subgroup will
move long lists of rows at low edit cost, while a simultaneous half-turn
will remain separated under every alignment. After proving separation, we derive a common analytic comparison and verify
its arithmetic hypothesis for these translations. The tags will then be
replaced by binary blocks in the next section.

\subsection{Translated rows on a finite circle}
Fix a sufficiently large integer $h$. Choose $h^2$ distinct primes
$p_{\ell k}\in[h^3,2h^3]$, for $1\leq\ell,k\leq h$. Here $\pi(x)$
denotes the number of primes at most $x$. The estimate
$\pi(2x)-\pi(x)>3x/(5\log x)$ for sufficiently large $x$
\cite[Corollary~3]{RosserSchoenfeld1962} supplies this many primes for every
sufficiently large $h$. Put
\begin{equation}\label{mask:parameters}
 P_\ell=\prod_{k=1}^h p_{\ell k},\qquad
 M=2\prod_{\ell=1}^hP_\ell,\qquad R=h^{2h},\qquad
 G=(M^{-1}\mathbb Z)/\mathbb Z.
\end{equation}
The factor two puts the half-turn in the finite grid $G$.

At level $\ell$ the child set is
\[
 J_\ell=\{(k,b):1\leq k\leq h,\ 0\leq b<p_{\ell k}\},
 \qquad t_\ell=\sum_kp_{\ell k},
\]
with a fixed order. For a residue vector
$z\in\prod_k\mathbb Z/p_{\ell k}\mathbb Z$, define the additive maps
\begin{equation}\label{mask:gamma}
 \sigma_\ell(z)=\sum_k\frac{z_k}{p_{\ell k}},\qquad
 \gamma_{\ell,(k,b)}(z)=\sigma_\ell(z)+\frac{bz_k}{p_{\ell k}}
 \quad\text{in }G.
\end{equation}
For an integer $a$, write $z(a)$ for its vector of residues at this level.
By the Chinese remainder theorem, $a\bmod P_\ell$ runs through every $z$.

Set $r_0=1$ and $r_\ell=t_\ell r_{\ell-1}$. A state
$v\in G^{r_\ell}$ is a tuple $(v_j)_{j\in J_\ell}$ with
$v_j\in G^{r_{\ell-1}}$. For $\theta\in G$, write $\theta\mathbf1$ for the
vector whose every coordinate is $\theta$. At level zero, $S_0(v)$ is the length-$M$ binary word with
bit at position $b\in\{0,\ldots,M-1\}$ equal to
\[
 \mathbf1_{[0,1/2)}(b/M-v),
\]
where the interval is read modulo one. Thus $S_0(v)$ records $M/2$
consecutive sites of the circle.

For $\ell\geq1$, define $S_\ell(v)$ by concatenating, for
$a=0,\ldots,RP_\ell-1$, the row
\begin{equation}\label{mask:word}
 \mathop{\Vert}_{j\in J_\ell}
 [j,S_{\ell-1}(v_j+\gamma_{\ell,j}(z(a))\mathbf1)].
\end{equation}
Here $\Vert$ means concatenation in the fixed order, and $[j,U]$ attaches
the tag $j$ to every letter of $U$. The tags are part of the alphabet, so a
match must preserve the child index. All level-$\ell$ words have length
\begin{equation}\label{mask:length}
 L_0=M,\qquad L_\ell=RP_\ell t_\ell L_{\ell-1}.
\end{equation}
At level $h$ the alphabet has size $A=2r_h$. Injectivity of the map from
states to words is not required.

The following subgroup records the translations that move row indices:
\begin{equation}\label{mask:subgroup}
 C_0=\{0\},\qquad
 C_\ell=C_{\ell-1}^{J_\ell}+
 \{(\gamma_{\ell,j}(z)\mathbf1)_{j\in J_\ell}:
        z\in\textstyle\prod_k\mathbb Z/p_{\ell k}\mathbb Z\}.
\end{equation}
For a leaf coordinate $i$, let $e_i(\theta)$ have value $\theta$ there
and zero in the other coordinates.

\begin{lemma}[Cheap translations]\label{mask:cheap}
For every $0\leq\ell\leq h$, every $v\in G^{r_\ell}$, and every
$c\in C_\ell$,
\begin{equation}\label{mask:cheap-subgroup}
 \ED(S_\ell(v),S_\ell(v+c))\leq\frac{2\ell L_\ell}{R}.
\end{equation}
At the top level, for every leaf $i$ and $1\leq a\leq M/2$,
\begin{equation}\label{mask:cheap-leaf}
 \ED(S_h(v),S_h(v+e_i(a/M)))\leq\frac{2aL_h}{Mr_h}.
\end{equation}
Both estimates hold at every base state $v$.
\end{lemma}
\begin{proof}
For a translation from the second summand of~\eqref{mask:subgroup}, choose
its row shift $q\in\{0,\ldots,P_\ell-1\}$ by the Chinese remainder
theorem. Additivity of~\eqref{mask:gamma} identifies the rows after this
shift. Deleting $q$ rows at one end and inserting $q$ at the other costs
at most $2L_\ell/R$. For a translation in the first summand, apply the
previous-level bound inside each child block of each row. Its total cost is
at most $2(\ell-1)L_\ell/R$. These bounds hold at every state, so the
triangle inequality combines them at the intermediate state and proves
\eqref{mask:cheap-subgroup} by induction.

A leaf occupies exactly the fraction $1/r_h$ of the root word, in copies of
$S_0$. Shifting its half-circle by $a/M$, for $a\leq M/2$, changes exactly
$2a$ bits in each copy. Substituting those bits, with every tag unchanged,
proves~\eqref{mask:cheap-leaf}.
\end{proof}

\subsection{A scalar cost and separation under arbitrary alignments}
The subgroup can choose different row translations at successive levels.
To measure how much of a uniform phase survives those choices, define
\begin{equation}\label{mask:phase-cost}
 D_0(\theta)=\|\theta\|_{\mathbb R/\mathbb Z},\qquad
 D_\ell(\theta)=\min_z\frac1{t_\ell}\sum_{j\in J_\ell}
        D_{\ell-1}(\theta+\gamma_{\ell,j}(z)).
\end{equation}
The minimum exists because its set is finite. These functions are symmetric
and subadditive: negate a minimizing residue vector for symmetry, and add
two minimizing vectors before using the preceding subadditivity inequality.
The zero vector also shows $0\leq D_\ell\leq1/2$.

\begin{lemma}[Word separation]\label{mask:separation}
For every $0\leq\ell\leq h$, $v\in G^{r_\ell}$, and $\theta\in G$,
\begin{equation}\label{mask:separation-bound}
 r_{L_\ell}(S_\ell(v),S_\ell(v+\theta\mathbf1))
 \geq\frac{2^{-\ell}}4L_\ell D_\ell(\theta).
\end{equation}
\end{lemma}
\begin{proof}
At level zero let $E$ be the deficit. In a matching of $M-E$ pairs, the
$k$th selected position on either side lies between $k$ and $k+E$, so each
matched pair has displacement at most $E$. The Hamming distance of the two
half-circle words is $2M\|\theta\|_{\mathbb R/\mathbb Z}$. A position where
their bits disagree is either unmatched in the source or matched across a
transition of the target word. The target has at most two transitions in
its linear ordering, and at most $4E$ positions lie within distance $E$ of
them. Thus the Hamming distance is at most $E+4E$. It follows that
$E\geq2MD_0(\theta)/5\geq MD_0(\theta)/4$.

Suppose the assertion holds at level $\ell-1$. Fix a longest common
subsequence matching of the level-$\ell$ words. Let $t=t_\ell$ and
$L=L_{\ell-1}$. If a source row has all its matches in a target row with
indices $a,a'$, respectively, its child $j$ can match only child $j$ in
that target row. Their phase difference is
$(\theta+\gamma_{\ell,j}(z(a'-a)))\mathbf1$. By induction, the number of
source positions lost in this row is at least
\[
 \frac{2^{-(\ell-1)}}4L\sum_j
 D_{\ell-1}(\theta+\gamma_{\ell,j}(z(a'-a)))
 \geq\frac{2^{-(\ell-1)}}4tL D_\ell(\theta).
\]
A row with no matches loses every source position and satisfies the same
bound. Hence Lemma~\ref{pre:tagged-rows} applies with
$A=2^{-(\ell-1)}D_\ell(\theta)/4\leq1/8$. Since $t\geq4$, its denominator
$1+tA/(t-1)$ is less than two. We obtain a deficit at least
$AL_\ell/2$, which is~\eqref{mask:separation-bound}.
\end{proof}

The remaining geometric question is now one-dimensional: how small can
$D_h(1/2)$ be? The next descent gives the quantitative answer. It is here
that the prime groups and their complete residue lists enter the proof.

\subsection{Keeping a nonzero phase through the descent}
\begin{lemma}[Denominator resets]\label{mask:reset}
For all sufficiently large $h$,
\begin{equation}\label{mask:reset-bound}
 D_h(1/2)\geq4^{-h}h^{-h/2-O(\sqrt h)-4h/100},
\end{equation}
with an absolute implicit constant.
\end{lemma}
\begin{proof}
At a level $\ell$, choose a minimizing residue vector $z$ in
\eqref{mask:phase-cost} and let $s=|\{k:z_k\ne0\}|$. If $z_k\ne0$,
multiplication by $z_k$ permutes the $p_{\ell k}$ grid points. For every
translate of that grid, symmetry and subadditivity give
\[
 D_{\ell-1}(1/p_{\ell k})
 \leq D_{\ell-1}(x+b/p_{\ell k})+
       D_{\ell-1}(x+(b+1)/p_{\ell k}).
\]
Averaging over $b$ says that the average cost in this child group is at
least $D_{\ell-1}(1/p_{\ell k})/2$. Its weight in the row is
$p_{\ell k}/t_\ell\geq1/(2h)$. Summing over active groups proves, when
$s>0$, the reset option
\begin{equation}\label{mask:reset-option}
 D_\ell(\theta)\geq\frac{s}{4h}
       \min_{k:z_k\ne0}D_{\ell-1}(1/p_{\ell k}).
\end{equation}
If $s<\sqrt h$, the inactive groups have total weight at least
$(h-s)/(2h)\geq1/4$ for large $h$. They all have phase
$\theta+\sigma_\ell(z)$, so there is also the option
\begin{equation}\label{mask:inactive-option}
 D_\ell(\theta)\geq\tfrac14D_{\ell-1}(\theta+\sigma_\ell(z)).
\end{equation}

Set $K=\lfloor h/100\rfloor$ and start at level $h$ with phase $1/2$.
At entry to level $\ell$, maintain a nonzero phase whose reduced denominator
is a product of at most $K$ distinct primes, all equal to two or belonging
to levels strictly above $\ell$. The initial phase has this property.

Use~\eqref{mask:reset-option} whenever $s\geq\sqrt h$. Use it also when
$s<\sqrt h$ but the current denominator count plus $s$ exceeds $K$.
In the latter case $s>0$, so the option is available. Continue at a prime
reciprocal attaining its minimum. A reset restores the count to one and
leaves a nonzero phase with a prime from the current level.

In every other case use~\eqref{mask:inactive-option}. If the current phase
is $a/Q$ in lowest terms, then $Q>1$. The new summand has denominator
dividing a product $P$ of active current-level primes, with $\gcd(Q,P)=1$.
If $a/Q+b/P$ were an integer, reduction of $aP+bQ$ modulo $Q$ would give
$aP\equiv0\pmod Q$, impossible. The resulting nonzero phase has at most
the old denominator count plus $s$ factors. Thus the invariant is preserved.

Only $O(\sqrt h)$ resets have $s<\sqrt h$. For each such reset, take the
interval of steps since the preceding reset of either kind, or since the
start. Its initial count is one. Overflow implies that the sum of its small
supports, including the last one, exceeds $K-1$. These intervals are
disjoint, and their total support sum is at most $h\sqrt h$. The number $q$
of small resets therefore satisfies
\[
 q\leq\frac{h\sqrt h}{K-1}=O(\sqrt h).
\]
This count remains valid when large resets occur between small ones.

A large reset costs a factor at least $1/(4\sqrt h)$, an inactive step at
least $1/4$, and a small reset at least $1/(4h)$. Their product is at least
$4^{-h}h^{-h/2-q/2}$. At level zero the nonzero phase has at most $K$
denominator primes, each at most $2h^3\leq h^4$. Its distance to zero is
therefore at least $h^{-4K}$. Multiplying the factors and using
$K\leq h/100$ gives~\eqref{mask:reset-bound}.
\end{proof}

Lemmas~\ref{mask:separation} and~\ref{mask:reset} give a pointwise half-turn
deficit at least $L_hh^{-(0.54+o(1))h}$. The exponential factors $8^{-h}$
are absorbed by the $o(1)h$ term on this logarithmic scale. We next check
that the cheap subgroup forces the large frequency size needed to compare
this separation with an $\ell_1$ image.

\subsection{A common analytic comparison for circle shifts}

The word construction has supplied cheap translations and an expensive
half-turn. We now derive the inequality that compares those costs in an
arbitrary $\ell_1$ image. The proof is stated for any finite circle product;
the next subsection will verify its arithmetic hypothesis for the translation
subgroup $C_h$.

For a finite set $X$ and $A\subset X$, write
$\delta_A(x,y)=|\mathbf1_A(x)-\mathbf1_A(y)|$. The finite cut
representation of $\ell_1$ distances is standard
\cite[Section~4]{LinialLondonRabinovich1995}. We include the summability
argument because an $\ell_1$ target may have infinitely many coordinates.

\begin{lemma}[Finite cut representation]\label{pre:cuts}
For $f:X\to\ell_1$ with $X$ finite, there are nonnegative coefficients
$a_A$ on the nonempty proper subsets of $X$ such that
\[
 \|f(x)-f(y)\|_1=\sum_A a_A\delta_A(x,y),\qquad
 \sum_Aa_A\leq\sum_{\{x,y\}\subset X}\|f(x)-f(y)\|_1.
\]
\end{lemma}
\begin{proof}
For real $s,t$,
$|s-t|=\int_{\mathbb R}|\mathbf1_{s>u}-\mathbf1_{t>u}|\,du$.
Apply this to each coordinate of $f$ and group threshold sets that define
the same subset of $X$. Empty and full cuts contribute zero. Every other
cut separates some unordered pair, so its weight is counted at least once
in the displayed pair sum, which is finite. Nonnegative sums and integrals
may be interchanged, giving the representation and the bound.
\end{proof}

For this common lemma, let $M\geq2$ be any even integer,
$G=(M^{-1}\mathbb Z)/\mathbb Z$, and $V=G^r$ for
$r\geq1$. For a map $F:V\to\ell_1$ define its averaged displacement by
\[
 \Delta_F(a)=\E_{v\in V}\|F(v+a)-F(v)\|_1.
\]
All averages on finite groups in this paper are uniform. Represent a
character of $V$ as
$\chi_\alpha(v)=\exp(2\pi\mathrm i\sum_i\alpha_i v_i)$, choosing each
integer $\alpha_i$ with $|\alpha_i|\leq M/2$. Either representative at an
endpoint gives the same absolute value. Put
$T(\alpha)=\sum_i|\alpha_i|$ and
$\mathcal D_M=\{1,2,4,\ldots\}\cap[1,M/2]$.
For $\theta\in G$, write $\theta e_i$ for the vector in $V$ with value
$\theta$ in coordinate $i$ and zero elsewhere.

\begin{lemma}[Weighted finite-circle inequality]\label{pre:weighted-circle}
Let $H\leq V$, let $w=(1/2,\ldots,1/2)\in V$, and suppose $B>0$ satisfies
\begin{equation}\label{pre:branching-premise}
 \chi_\alpha|_H=1,\quad \chi_\alpha(w)\ne1
 \quad\Longrightarrow\quad T(\alpha)\geq B.
\end{equation}
For every $F:V\to\ell_1$ there are nonnegative weights $\omega_\alpha$ with
\begin{equation}\label{pre:fourier}
 \Delta_F(a)=\sum_\alpha\omega_\alpha|1-\chi_\alpha(a)|^2
 \qquad(a\in V).
\end{equation}
Writing $H^\perp=\{\alpha:\chi_\alpha|_H=1\}$, these weights satisfy
\begin{align}
 \sum_{\alpha\notin H^\perp}\omega_\alpha
   &=\tfrac12\E_{u\in H}\Delta_F(u),\label{pre:off-mass}\\
 4\sum_\alpha\omega_\alpha T(\alpha)
   &\leq\sum_{i=1}^r\sum_{a\in\mathcal D_M}
        \frac Ma\Delta_F((a/M)e_i).\label{pre:weighted-mass}
\end{align}
Consequently
\begin{equation}\label{pre:displacement}
 \Delta_F(w)\leq2\E_{u\in H}\Delta_F(u)
 +\frac1B\sum_{i=1}^r\sum_{a\in\mathcal D_M}
        \frac Ma\Delta_F((a/M)e_i).
\end{equation}
\end{lemma}
\begin{proof}
For a scalar function $g$ on $V$, normalized Fourier coefficients and
Parseval give
\[
 \E_v|g(v+a)-g(v)|^2
 =\sum_\alpha|\widehat g(\alpha)|^2|1-\chi_\alpha(a)|^2.
\]
For $g=\mathbf1_A$ its squared difference is the cut distance. Multiply
by the coefficients of Lemma~\ref{pre:cuts} and sum over cuts. This gives
\eqref{pre:fourier} with
$\omega_\alpha=\sum_Aa_A|\widehat{\mathbf1_A}(\alpha)|^2\geq0$.
The sums are finite after grouping cuts. Since a nontrivial character has
mean zero on a finite subgroup, the mean of $|1-\chi_\alpha(u)|^2$ on $H$
is two off $H^\perp$ and zero on it. This proves~\eqref{pre:off-mass}.

For each nonzero coordinate, take the largest power of two
$a\leq M/(2|\alpha_i|)$. It belongs to $\mathcal D_M$ and obeys
\[
 \frac14<\frac{a|\alpha_i|}{M}\leq\frac12,
 \qquad |1-e^{2\pi\mathrm i a\alpha_i/M}|^2\geq2,
 \qquad \frac Ma\geq2|\alpha_i|.
\]
That single scale contributes at least $4|\alpha_i|$. Sum over coordinates
and then against the nonnegative Fourier weights to obtain
\eqref{pre:weighted-mass}.

The half-turn multiplier is zero or four. Its detecting frequencies off
$H^\perp$ contribute at most $4\sum_{\alpha\notin H^\perp}\omega_\alpha$.
By~\eqref{pre:branching-premise}, the remaining ones contribute at most
$(4/B)\sum_\alpha\omega_\alpha T(\alpha)$. The two preceding bounds give
\eqref{pre:displacement}. Equivalently, the same estimates prove that
inequality for each character multiplier separately and then sum it with
nonnegative weights. Thus the lemma supplies both the weighted-mass and
the termwise-displacement forms of the comparison.
\end{proof}

The hypothesis~\eqref{pre:branching-premise} is the only arithmetic input
to this lemma. We shall prove it from the prime translations of each
construction. The inequalities themselves concern averages over all base
states; they impose no pointwise condition on individual Fourier weights.


\subsection{Prime groups force leaf frequency}
Return to the value of $M$ and the grid $G$ in~\eqref{mask:parameters}.
Put $r=r_h$, $V=G^r$, $w=\tfrac12\mathbf1$, and $B=(h/2)^h$. For a character
$\chi_\alpha$ choose the centered integer representatives used in
Lemma~\ref{pre:weighted-circle}.

\begin{lemma}[Frequency branching]\label{mask:branching}
If $\chi_\alpha|_{C_h}=1$ and $\chi_\alpha(w)\ne1$, then
\begin{equation}\label{mask:weight}
 \sum_{i=1}^r|\alpha_i|\geq B.
\end{equation}
\end{lemma}
\begin{proof}
View the state coordinates as the leaves of the tree in~\eqref{mask:word}.
At a node let $\beta$ be the sum of the integer leaf frequencies below it,
and at a level-$\ell$ node let $\beta_{(k,b)}$ be its child sums.
At this node, the translation defined by a residue vector through the maps
$\gamma_{\ell,j}$, and set to zero outside the node, belongs to $C_h$ by the
independent child summands in~\eqref{mask:subgroup}. Choosing $z_k=1$ and
all other coordinates zero gives the congruences
\begin{equation}\label{mask:congruence}
 \beta+\sum_{b=0}^{p_{\ell k}-1}b\beta_{(k,b)}
       \equiv0\pmod{p_{\ell k}}\qquad(1\leq k\leq h).
\end{equation}
They are unchanged if a frequency representative is altered by $M$.

Suppose the sum in~\eqref{mask:weight} were less than $B$. Every nonzero
node sum would have absolute value less than $B$. If $q$ distinct
current-level primes divide such a sum, then
\[
 h^{3q}\leq|\beta|<(h/2)^h<h^h,
\]
so $q<h/3$. For every other prime group,~\eqref{mask:congruence} forces a
nonzero child sum in that group. Each nonzero internal node therefore has
at least $h/2$ nonzero children in distinct groups. At the root, the sum is
odd because $\chi_\alpha(w)\ne1$; parity is well defined since $M$ is even.
Iteration gives at least $(h/2)^h$ nonzero integer leaf frequencies,
contradicting the assumed total absolute sum.
\end{proof}

The lemma verifies~\eqref{pre:branching-premise} with $H=C_h$ and
$B=(h/2)^h$. Thus every $F:V\to\ell_1$ satisfies
\begin{equation}\label{mask:fourier-bound}
 \Delta_F(w)\leq2\E_{c\in C_h}\Delta_F(c)+
 \frac1B\sum_{i=1}^r\sum_{a\in\mathcal D_M}\frac Ma
       \Delta_F(e_i(a/M)).
\end{equation}
We now have all three ingredients: the pointwise half-turn separation,
the pointwise cheap translations, and the averaged $\ell_1$ comparison.
The tags still use $2r_h$ letters. The next section supplies the binary
substitution, after which Section~\ref{sec:mask-binary} completes the
numerical distortion estimate.
